\subsection{\texorpdfstring{\(R^{n}\) 中的仿射线性簇}{R 的 n 次幂中的仿射线性簇}}

对给定的 \(p\) 维仿射线性簇 \(\mathbf{V}\) （\(\mathbf{R}^n\) 中的），存在一个仿射映射 \(f\) （\(\mathbf{R}^n\) 到其自身的），它把 \(\mathbf{V}\) 变换成一个 \(p\) 维坐标簇，也就是说，变换成一个向量子空间 \(\mathbf{V}'\) ——由 \(p\) 个典范基 \((\boldsymbol{e}_i)\) （\(\mathbf{R}^n\) 的）中的向量所生成。存在 \(\mathbf{V}'\) 到 \(\mathbf{R}^p\) 上的一个映射，它是 \(\mathbf{V}'\) 的向量空间结构和拓扑到 \(\mathbf{R}^p\) 的相应结构上的一个同构（此外，把 \(\mathbf{R}^p\) 与这样的坐标簇 \(\mathbf{V}'\) 等同起来常常是方便的，例如与由 \(\boldsymbol{e}_1, \boldsymbol{e}_2, \ldots, \boldsymbol{e}_p\) 生成的向量子空间等同）。此外，\(\mathbf{V}'\) 是 \(\mathbf{R}^n\) 的闭子集（第 I 章，§ 4，第 3 小节，命题 7 的推论）；因此：

命题 2. \(\mathbf{R}^n\) 中每个 p 维仿射线性簇是 \(\mathbf{R}^n\) 的闭子集，且同胚于 \(\mathbf{R}^p\)。

正是这个结果产生了在任意除环上的向量空间中把一维与二维的仿射线性簇称为直线与平面这些名称的由来。我们还提醒，对 \(n \geqslant 1\) ，\(R^{n}\) 的 n - 1 维仿射线性簇称为超平面（出处同上）。

n 个一维坐标簇，也就是说，分别过 o 与 n 个点 \(e_{i}\) 的 n 条直线，称为坐标轴。当 n = 2 时，过 \(e_{1}\) 的轴有时称为横轴，过 \(e_{2}\) 的轴称为纵轴；点 \(x \in R^{2}\) 的第一个坐标称为它的横坐标，第二个坐标称为它的纵坐标。

过点 \(\pmb{a}\) 的每条直线 D 都有参数表示 \(t \to \pmb{a} + tb\) ，其中 \(t\) 取遍 \(\mathbf{R}\) ，而 \(b \neq 0\) ；这个映射是 \(\mathbf{R}\) 到 \(D\) 上的一个同胚。向量 \(\pmb{b}\) 称为 D 的方向向量，它的分量 \(b_{i}\) （ \(1 \leqslant i \leqslant n\) ）称为 D 的方向比。若 \(b'\) 是 D 的另一个方向向量，则存在 R 中 \(h \neq o\) 使 \(b' = hb\)。

点 \(a + tb\) （t 取遍 \(\geqslant 0\) 的实数的集）组成的集称为以 a 为端点、以 b 为方向向量的闭射线（或简称射线，或半直线）（或以方向比 \(b_{i}\) 给出的）。它是 \(R^{n}\) 的闭子集，同胚于 R 的区间 \([0, +\infty[\) ，因此是连通的。直线 D 是以 a 为端点、分别以 b 与 -b 为方向向量的两条射线的并；这两条射线称为互为反向的射线。

滥用地说，点 \(\pmb{a} + t\pmb{b}\) （\(t\) 取遍 \(> 0\) 的实数的集）组成的集称为以 \(\pmb{a}\) 为端点、以 \(\pmb{b}\) 为方向向量的开射线；它同胚于区间 \(]0, +\infty[\) （从而也同胚于 \(\mathbf{R}\)），但它在 \(\mathbf{R}^n\) 中不是开的（当 \(n > 1\) 时），尽管它在包含它的直线中是开的。

过两个不同点 x 与 y 的直线也有参数表示 \((u, v) \to ux + vy\) ，其中 \((u, v)\) 取遍满足 \(u + v = 1\) 的实数对的集。对任意两点 x, y（不同与否），点 \(ux + vy\) ——其中 \((u, v)\) 取遍满足 \(u \geqslant 0\) 、\(v \geqslant 0\) 且 \(u + v = 1\) 的实数对的集——组成的集，称为以 x, y 为端点的闭线段（或简称线段）。闭线段是紧的且连通的，因为若它的端点不同，它同胚于 R 的区间{[}0, 1{]}。

若 \(x \neq y\) ，则点 \(ux + vy\) ——满足 \(u > 0, v > 0\) 且 \(u + v = 1\) ——组成的集（滥用地说）称为以 \(x, y\) 为端点的开线段；它同胚于开区间 \(]0, 1[\) （从而也同胚于 \(R\)）。最后，\(\{y\}\) 与以 \(x, y\) 为端点的开线段的并有时称为在 \(x\) 处开、在 \(y\) 处闭的线段；它同胚于区间 \([0, 1[\)。以 \(x\) 与 \(y\) 为端点的所有线段都是连通的，而且它们中每一个的闭包都是以同样端点为端点的闭线段。

命题 3. 设 \(f(x) = f(x_1, x_2, \ldots, x_n)\) 是定义在 \(\mathbf{R}^n\) 上的实系数多项式，不恒等于零。则集 \(\overline{f}(0)\) 的补集在 \(\mathbf{R}^n\) 中稠密。

设 \(\pmb{x}\) 是 \(\mathbb{R}^n\) 中任意一点，\(y \in \mathbb{R}^n\) 满足 \(f(y) \neq 0\) ；\(\varphi(t) = f(x + t(y - x))\) 是实变量 \(t\) 的多项式，不恒等于零；因此存在任意小的 \(t\) 值使 \(\varphi(t) \neq 0\) 。这表明 \(\pmb{x}\) 属于 \(\overline{f}(0)\) 的补集的闭包。

推论. 维数 \(p < n\) 的仿射线性簇的补集在 \(\mathbf{R}^n\) 中稠密。

由于每个维数 p \textless{} n 的仿射线性簇都含于一个超平面中，只需对超平面证明该推论；而超平面由方程 \(g(x) = 0\) 定义，其中 g 是一个不恒等于零的线性多项式。

命题 4. 在 \(\mathbf{R}^n\) （ \(n \geqslant 1\) ）中，每个超平面的补集都有两个连通分支。

设 \(g(x) = 0\) 是超平面 \(H\) （\(\mathbb{R}^n\) 中的）的一个方程，其中 \(g\) 是线性多项式。集 \(\complement H\) 是集 \(E_1\) （由一切点 \(x\) ——满足 \(g(x) > 0\) 的——组成）与集 \(E_2\) （由一切点 \(x\) ——满足 \(g(x) < 0\) 的——组成）的并。\(E_1\) 与 \(E_2\) 都是连通的，因为若 \(g(x) > 0\) 且 \(g(y) > 0\) ，则我们有 \(g(ux + vy) = ug(x) + vg(y) > 0\) 只要 \(u \geqslant 0\) 、\(v \geqslant 0\) 且 \(u + v = 1\) ；换句话说，以 \(x\) 与 \(y\) 为端点的线段含于 \(E_1\) 。对 \(E_2\) 同理。另一方面，\(\complement H\) 不是连通的，因为它在 \(R\) 中的、经 \(g\) 而得的象是区间 \(]0, +\infty[ \text{与}] - \infty, o[\) 的并。

超平面 H 的补集的分支 \(E_{1}\) 、\(E_{2}\) 称为由 H 决定的开半空间。

\(E_{1}\) 与 \(E_{2}\) 的闭包，即分别为 \(E_{1} \cup H\) 与 \(E_{2} \cup H\) ，称为由 H 决定的闭半空间。

注意，\(R^{n}\) 到其自身的、把 H 变换成“坐标”超平面（例如方程为 \(x_{n}=0\) 的超平面）的仿射映射，也把由 H 决定的开半空间分别变换成由关系 \(x_{n}>0\) 与 \(x_{n}<0\) 所定义的开半空间；后者是开长方体，因此同胚于 \(R^{n}\)。

